\section{Conclusion}\label{sec:conclusion}
We asked whether explanations of ML-based intrusion detectors in cyber-physical systems reach operators correctly and in time, and argued that XAI for CPS should be built on the control structure and evaluated against the physical deadlines of the process. With deadline-conditioned metrics, cross-checked loop-level ground truth for HAI~21.03 and a closed-loop simulator with hazard deadlines, we found that score attributions of every kind degrade as attacks propagate through coupled loops; that temporal accumulation recovers part of this loss for any explainer; that grouping by control loops is essential; and that a single physical prior, separating the persistent departure of the manipulated loop from the oscillatory compensation of the others, makes a physics-informed explainer significantly more accurate than all score attributions, both at the alarm and a minute later. Most remaining deadline failures are detection failures.

Future work should replicate the study on SWaT, WADI, newer HAI releases and other CPS domains; test physics-informed explainers against adaptive adversaries who steer the explanation while keeping the alarm; measure latency on embedded hardware; and evaluate deadline-aware explanations with operators.
